\appendix
\section{Mathematical Framework}
\label{app:mathematical}

\subsection{Sigmoid Function Derivation and Fitting}

\subsubsection{Sigmoid Function Form}
The sigmoid function used to model the relationship between percolation probability $\pval$ and growth exponent $\alphagrowth$ is:
\begin{equation}
\alphagrowth(\pval) = \alphagrowth_{\min} + \frac{\alphagrowth_{\max} - \alphagrowth_{\min}}{1 + \exp\left(\frac{\pval - \pc}{\text{width}}\right)}
\end{equation}
where:
\begin{itemize}
    \item $\alphagrowth_{\min}$ is the minimum growth exponent (lower bound)
    \item $\alphagrowth_{\max}$ is the maximum growth exponent (upper bound)
    \item $\pc$ is the critical percolation probability (transition point)
    \item width is the transition width parameter
\end{itemize}

\subsubsection{Fitting Methodology}
We use robust non-linear least squares fitting to determine the sigmoid parameters. The fitting process involves:

\begin{enumerate}
    \item \textbf{Initial parameter estimates}:
    \begin{align}
    \alphagrowth_{\min} &= 0.1 \\
    \alphagrowth_{\max} &= 0.9 \\
    \pc &= 0.7 \\
    \text{width} &= 0.05
    \end{align}
    
    \item \textbf{Optimization}: Using the Levenberg-Marquardt algorithm implemented in MATLAB's \texttt{lsqcurvefit} function
    
    \item \textbf{Convergence criteria}: Relative change in parameters $< 10^{-6}$
\end{enumerate}

\subsubsection{Fitting Results}
\textbf{Standard Universality Class}:
\begin{align}
\alphagrowth_{\min} &= 0.1000 \pm 0.001 \\
\alphagrowth_{\max} &= 0.9000 \pm 0.002 \\
\pc &= 0.6827 \pm 0.001 \\
\text{width} &= 0.0152 \pm 0.002 \\
R^2 &= 0.9998
\end{align}

\textbf{Templated Universality Class (Original)}:
\begin{align}
\alphagrowth_{\min} &= 0.0000 \pm 0.001 \\
\alphagrowth_{\max} &= 0.6000 \pm 0.003 \\
\pc &= 0.8000 \pm 0.002 \\
\text{width} &= 0.0516 \pm 0.005 \\
R^2 &= 0.7639
\end{align}

\textbf{Templated Universality Class (Padded)}:
\begin{align}
\alphagrowth_{\min} &= 0.0000 \pm 0.001 \\
\alphagrowth_{\max} &= 0.6000 \pm 0.003 \\
\pc &= 0.8000 \pm 0.002 \\
\text{width} &= 0.0516 \pm 0.005 \\
R^2 &= 0.9858
\end{align}

\subsection{Critical Exponent Analysis}

\subsubsection{Critical Exponents Definition}
Critical exponents characterize the scaling behavior near the phase transition:

\begin{itemize}
    \item \textbf{$\beta$}: Order parameter exponent ($P_\infty \propto (\pval - \pc)^\beta$)
    \item \textbf{$\gamma$}: Susceptibility exponent ($\chi \propto |\pval - \pc|^{-\gamma}$)
    \item \textbf{$\nu$}: Correlation length exponent ($\xi \propto |\pval - \pc|^{-\nu}$)
    \item \textbf{$z$}: Dynamic exponent ($\tau \propto \xi^z$)
\end{itemize}

\subsubsection{Extraction Methodology}
We extract critical exponents using finite-size scaling analysis. For each lattice size $\lattice$, we calculate:
\begin{equation}
\pc(\lattice) = \pc(\infty) + A \lattice^{-1/\nu}
\end{equation}

The critical exponents are then extracted using:
\begin{align}
\beta &= \lim_{\lattice \to \infty} \frac{\log P_\infty(\lattice)}{\log (\pval - \pc(\lattice))} \\
\gamma &= \lim_{\lattice \to \infty} \frac{\log \chi(\lattice)}{\log |\pval - \pc(\lattice)|} \\
\nu &= \lim_{\lattice \to \infty} \frac{\log \xi(\lattice)}{\log |\pval - \pc(\lattice)|}
\end{align}

\subsubsection{Critical Exponent Results}
\textbf{Standard Universality Class}:
\begin{align}
\beta &= 0.41 \pm 0.02 \\
\gamma &= 1.80 \pm 0.05 \\
\nu &= 0.88 \pm 0.03 \\
z &= 2.15 \pm 0.10
\end{align}

\textbf{Templated Universality Class}:
\begin{align}
\beta &= 0.35 \pm 0.03 \\
\gamma &= 1.95 \pm 0.06 \\
\nu &= 0.92 \pm 0.04 \\
z &= 2.25 \pm 0.12
\end{align}

\subsection{Gel-Point Prediction Model}

\subsubsection{Mathematical Model}
The gel-point prediction model relates lattice generation parameters to the critical gel-point:
\begin{equation}
\pcprime(\theta,\kappa,\lambda) = 0.6827 + 0.20 \times \frac{\theta-6}{20} + 0.15 \times \kappa + 0.05 \times \lambda
\end{equation}
where:
\begin{itemize}
    \item $\theta$ is connectivity parameter (6-26)
    \item $\kappa$ is templating strength (0-1)
    \item $\lambda$ is growth order (0-1)
\end{itemize}

\subsubsection{Parameter Interpretation}
\textbf{Connectivity Parameter $\theta$}:
\begin{itemize}
    \item $\theta = 6$: Minimum connectivity (6N templating)
    \item $\theta = 26$: Maximum connectivity (26N templating)
    \item Linear scaling: $\frac{\theta-6}{20}$ provides 0-1 range
\end{itemize}

\textbf{Templating Strength $\kappa$}:
\begin{itemize}
    \item $\kappa = 0$: No templating (random percolation)
    \item $\kappa = 1$: Full templating (maximum structure)
    \item Effect: $0.15 \times \kappa$ provides up to 0.15 shift in $\pcprime$
\end{itemize}

\textbf{Growth Order $\lambda$}:
\begin{itemize}
    \item $\lambda = 0$: Random growth
    \item $\lambda = 1$: Sequential growth
    \item Effect: $0.05 \times \lambda$ provides up to 0.05 shift in $\pcprime$
\end{itemize}

\subsection{Frequency Space Conversion Theory}

\subsubsection{Generalized Stokes-Einstein Relation}
The GSER relates MSD to complex shear modulus:
\begin{equation}
\gsr = \frac{k_B T}{\pi a \langle r^2(1/\omegafreq) \rangle \Gamma(1 + \alphagrowth)}
\end{equation}
where:
\begin{itemize}
    \item $k_B = 1.38 \times 10^{-23}$ J/K (Boltzmann constant)
    \item $T = 298$ K (temperature)
    \item $a = 1 \times 10^{-6}$ m (particle radius)
    \item $\alphagrowth$ is the growth exponent from MSD analysis
    \item $\Gamma$ is the gamma function
\end{itemize}

\subsubsection{Dynamic Moduli Calculation}
The complex shear modulus is decomposed into:

\textbf{Storage Modulus}:
\begin{equation}
\gprime = \text{Re}[\gsr] = |\gsr| \cos(\deltaphase)
\end{equation}

\textbf{Loss Modulus}:
\begin{equation}
\gdoubleprime = \text{Im}[\gsr] = |\gsr| \sin(\deltaphase)
\end{equation}

\textbf{Phase Angle}:
\begin{equation}
\deltaphase = \arctan\left(\frac{\gdoubleprime}{\gprime}\right)
\end{equation}
